\RequirePackage[table]{xcolor}
\documentclass[fonts]{icst}

\usepackage{moreverb}
\usepackage[breaklinks,colorlinks,bookmarksopen,bookmarksnumbered,linkcolor=ICSTblue,citecolor=blue,urlcolor=ICSTblue]{hyperref}
\usepackage{iftex}
\ifPDFTeX\usepackage{breakurl}\fi
\usepackage{doi}
\usepackage{xurl}
\usepackage[utf8]{inputenc}
\usepackage[english]{babel}
\usepackage{amsmath,amssymb,amsthm}
\usepackage{graphicx,flushend}
\usepackage{booktabs}
\usepackage{array}
\usepackage{multirow}
\usepackage{float}
\usepackage{enumitem}
\usepackage{dblfloatfix}
\usepackage[hypcap=false]{caption}
\usepackage{adjustbox}
\usepackage{balance}
\usepackage{tikz}
\usetikzlibrary{positioning,calc}
% float placement for a two-column layout: let floats sit near their first mention
\renewcommand{\topfraction}{0.85}\renewcommand{\dbltopfraction}{0.85}
\renewcommand{\bottomfraction}{0.5}\renewcommand{\textfraction}{0.1}
\renewcommand{\floatpagefraction}{0.75}\renewcommand{\dblfloatpagefraction}{0.75}
\setcounter{topnumber}{3}\setcounter{bottomnumber}{2}\setcounter{totalnumber}{5}\setcounter{dbltopnumber}{2}

% numbers shared with the BiasLN revision (tools/paper_assets.py)
\IfFileExists{generated/numbers.tex}{\input{generated/numbers.tex}}{}
% revised passages in blue for the editor's marked copy (build.sh sets revmarktrue)
\newif\ifrevmark\revmarkfalse
\ifrevmark\newcommand{\rev}[1]{\textcolor{blue}{#1}}\else\newcommand{\rev}[1]{#1}\fi
\newcommand{\num}[1]{\ifcsname #1\endcsname\csname #1\endcsname\else\textcolor{red}{??}\fi}
\newtheorem{proposition}{Proposition}
% \todo{...} marks text the authors must still check; remove all before submission.
\newcommand{\todo}[1]{\textcolor{red}{[TODO: #1]}}

\journalname{XXXXXX}
\articletype{Research Article}
\setcounter{page}{01}
\def\volumeyear{XXXX}
\def\volumenumber{XX}
\def\issuemonth{XX-XX}
\def\issuenumber{X}
\def\articleid{XX}
\def\copyrightholder{Xuan-Huy Dinh\textit{ et al.}, licensed to ICST}
\copyrightnote{This is an open access article distributed under the terms of the Creative Commons Attribution license (\url{http://creativecommons.org/licenses/by/3.0/}), which permits unlimited use, distribution and reproduction in any medium so long as the original work is properly cited.}
\def\articledoi{10.4108/XX.X.X.XX}
\received{XXXX}
\accepted{XXXX}
\published{XXXX}

\begin{document}
\runningheads{}{}

\title{A Survey of Parameter-Efficient Fine-Tuning for Generalizable Deepfake Detection}

\author{Xuan-Huy Dinh,\ Minh-Duc Vu,\ Thanh-Quan Phung-Le,\ Duc-Thinh Nguyen,\ Van-Dung Hoang\fnoteref{1}}

\address{
    Faculty of Information Technology\\
    Ho Chi Minh City University of Technology and Engineering (HCMUTE)\\
    Ho Chi Minh City, Vietnam
}

\abstract{
Detectors built on vision foundation models such as CLIP generalize well across deepfake benchmarks, and most of them adapt the frozen encoder with parameter-efficient fine-tuning (PEFT). This paper surveys PEFT for face deepfake detection. We group existing detectors into selective, additive, reparameterization and hybrid methods, discuss them along three research questions (how much to adapt, what and where to adapt, and how detectors transfer to unseen generators), and show that several common evaluation protocols select the reported checkpoint on test-benchmark data. We answer the questions with a controlled benchmark of eight schemes on CLIP ViT-L/14, re-trained with one pipeline, up to ten seeds and checkpoints selected only on validation data drawn from the same source as the training data. The selection rule alone moves one model's cross-dataset AUC from 0.934 to 0.923. Light adaptation is essential but sufficient: linear probing and full fine-tuning fall far behind, and larger budgets such as LoRA do not help. What is adapted matters more than how much: tuning only the LayerNorm parameters (0.1M), without any auxiliary loss, matches the strongest hybrid scheme and transfers better to unseen face-swapping methods than tuning biases, which instead raise in-domain accuracy. We close with recommendations and future directions.
}

\keywords{Deepfake detection, parameter-efficient fine-tuning, vision foundation models, CLIP, LayerNorm tuning, checkpoint selection, benchmark}

\fnotetext[1]{Corresponding author. Email: \email{dunghv@hcmute.edu.vn}}

\maketitle

% =====================================================================
\section{Introduction}

Face manipulation with GANs~\cite{goodfellow2020gan} and diffusion models~\cite{yang2022diffusion} has made realistic forged videos easy to produce, and detecting them is now a standard problem in media forensics~\cite{pei2024survey,liu2024evolving}. The central difficulty is generalization: a detector trained on one dataset, typically FaceForensics++ (FF++)~\cite{rossler2019ff}, must recognize forgeries produced by other pipelines, under other compression and capture conditions. Detectors that rely on low-level, method-specific artifacts degrade sharply on such shifts~\cite{qian2020frequency,luo2021srm,yan2023deepfakebench}.

Large pre-trained vision encoders changed this picture. Features of CLIP~\cite{radford2021clip_vitl} separate real from generated images surprisingly well even without training~\cite{ojha2023universal,tsai2024trainingfree}, and the best cross-dataset results are now obtained by adapting such an encoder rather than training a detector from scratch~\cite{yan2025effort,yermakov2025gend,cui2025forada}. Adapting it fully, however, is harmful: the deepfake training sets are small, and full fine-tuning distorts the pre-trained features~\cite{kumar2022finetuning} and fits the artifacts of the training forgeries. Parameter-efficient fine-tuning (PEFT) addresses both problems at once. It freezes almost all of the encoder, which preserves its general representation, and trains only a small set of existing parameters, added modules, prompts or low-rank updates, which also lowers training memory and storage.

PEFT detectors have multiplied quickly, but comparing them is hard. Each paper reports numbers from its own training recipe, test-set average and checkpoint-selection rule, and several of these rules pick the reported checkpoint on data from the test benchmarks instead of data from the training distribution (Sec.~\ref{sec:eval}). Differences of a few tenths of a point, often from a single run, are then read as methodological progress.

This paper surveys PEFT for face deepfake detection around three research questions (RQs):
\begin{itemize}[leftmargin=*, topsep=2pt, itemsep=2pt]
    \item \textbf{RQ1 (adaptation vs.\ preservation).} Why do schemes that adapt fewer parameters often reach higher cross-dataset accuracy than full fine-tuning, and where does the trade-off between fitting the training forgeries and preserving the pre-trained knowledge lie?
    \item \textbf{RQ2 (where and what to adapt).} Which parameters, layers and depths of a vision transformer should be adapted for deepfake detection?
    \item \textbf{RQ3 (unseen generators).} How well do PEFT detectors trained on one family of forgeries transfer to new generators?
\end{itemize}
Our contributions are:
\begin{enumerate}
    \item \textbf{A taxonomy} of PEFT-based deepfake detectors in four branches, selective, additive, reparameterization and hybrid, with their mechanisms, representative methods and trade-offs (Sec.~\ref{sec:taxonomy}).
    \item \textbf{A review of evaluation practice}, showing that checkpoint selection on test-benchmark data, inconsistent test-set averages and single-seed reporting make published numbers incomparable (Sec.~\ref{sec:eval}).
    \item \textbf{A controlled benchmark} that answers the three RQs on equal terms: eight schemes on CLIP ViT-L/14 and more than fifty configurations, re-trained with one pipeline and selected on FF++ validation data only, with up to ten seeds, paired significance and equivalence tests, per-dataset standard deviations, a parameter--accuracy trade-off analysis and measured costs (Sec.~\ref{sec:bench}).
    \item \textbf{An analysis} of the per-channel affine schemes (bias and LayerNorm tuning) explaining why they differ on unseen manipulations (Sec.~\ref{sec:analysis}), followed by recommendations and future directions (Secs.~\ref{sec:recommend} and~\ref{sec:open}).
\end{enumerate}
Code, configurations and all per-run results are available at \href{https://github.com/ducthinh4477/BiasLN}{https://github.com/ducthinh4477/BiasLN}.

% =====================================================================
\section{Background}
\label{sec:background}

\subsection{Problem setting and benchmarks}
Most face deepfake detectors operate on face crops extracted from video frames and output a fake probability per frame; a video score is the mean over frames. The common training source is FF++ (c23)~\cite{rossler2019ff}, which contains four manipulations (Deepfakes, Face2Face, FaceSwap, NeuralTextures). Generalization is measured \emph{cross-dataset}, on Celeb-DF-v2 (CDF-v2)~\cite{li2020celeb}, DeepFakeDetection (DFD)~\cite{dufour2019dfd}, the DFDC preview (DFDCP)~\cite{dolhansky2019dfdcp} and DFDC~\cite{dolhansky2020dfdc}, and \emph{cross-method}, on forgeries produced by newer generators applied to the same identities, as collected in DF40~\cite{yan2024df40}. DeepfakeBench~\cite{yan2023deepfakebench} standardized pre-processing (face detection, alignment, $1.3\times$ crop) and provides re-implementations of many detectors. Performance is reported as frame- or video-level AUC.

\subsection{Vision foundation models}
CLIP~\cite{radford2021clip_vitl} trains an image and a text transformer contrastively on 400M image--text pairs; its ViT-L/14 image encoder (303M parameters, 24 blocks) is the most common backbone in recent detectors. Self-supervised encoders such as DINO~\cite{caron2021emerging} are also used~\cite{zhang2026dino,zou2026clipdino}. Each transformer block applies a pre-norm LayerNorm, multi-head self-attention with query, key, value and output projections, a second LayerNorm and a two-layer MLP (fc1, fc2); all linear layers have biases.

\subsection{Parameter-efficient fine-tuning}
PEFT methods adapt a frozen pre-trained network with few trainable parameters. \emph{Selective} methods train an existing subset, such as all biases (BitFit~\cite{benzaken2022bitfit}) or the LayerNorm affine parameters~\cite{wang2022ln,giannou2023layernorm,valizadeh2024layernorm}. \emph{Reparameterization} methods add a low-rank update to frozen weights (LoRA~\cite{hu2022lora}) or tune part of a decomposition. \emph{Additive} methods insert small modules (adapters) or learnable input tokens (prompts). Selective and merged reparameterized updates leave the inference cost unchanged; unmerged reparameterizations and additive modules do not.

% =====================================================================
\section{Taxonomy of PEFT-Based Deepfake Detectors}
\label{sec:taxonomy}

We group PEFT-based detectors by \emph{how} they adapt the frozen encoder, which determines the number of trainable parameters, the inference cost and, as Sec.~\ref{sec:bench} shows, much of their behaviour. Fig.~\ref{fig:overview} shows where each branch acts in a transformer block, and Table~\ref{tab:taxonomy} summarizes the four branches; linear probing and full fine-tuning bound them from below and above.

\begin{figure*}[!t]
\centering
\includegraphics[width=0.95\textwidth]{figs/fig_overview_figma.pdf}
\caption{Where the four PEFT branches act in a frozen CLIP ViT-L/14 block (blue: frozen). Orange tags mark selective parameters ($\mathbf{b}$, $\boldsymbol\gamma$, $\boldsymbol\beta$), violet tags reparameterized updates ($\Delta\mathbf{W}$), green boxes additive modules, and the blue callout hybrid mechanisms.}
\label{fig:overview}
\end{figure*}

\begin{table*}[!t]
\centering
\caption{Taxonomy of PEFT schemes for deepfake detection. Infer.: effect on the inference cost of the frozen encoder; $\ast$: re-trained in our benchmark.}
\label{tab:taxonomy}
\footnotesize
\begin{tabular}{p{0.11\textwidth} p{0.20\textwidth} p{0.26\textwidth} p{0.07\textwidth} p{0.24\textwidth}}
\toprule
\textbf{Branch} & \textbf{Core mechanism} & \textbf{Representative methods} & \textbf{Infer.} & \textbf{Strengths / weaknesses} \\
\midrule
Selective & Freeze the encoder and unfreeze a small subset of its existing parameters (head, biases, LayerNorm, chosen blocks) & Linear probe / UnivFD~\cite{ojha2023universal}$^\ast$, training-free scoring~\cite{tsai2024trainingfree}, patch-token probes~\cite{abdullah2026tap}, BitFit~\cite{benzaken2022bitfit}$^\ast$, LayerNorm tuning~\cite{yermakov2025clip}$^\ast$, BiasLN$^\ast$, BPL~\cite{xu2026bpl}, layer masking~\cite{zhang2026layermask} & same & Lightest; architecture unchanged; strong preservation. Expressiveness limited to per-channel offsets and scales. \\
Additive & Insert small trainable modules (adapters, side networks) or learnable tokens (visual or text prompts) & Forensics Adapter~\cite{cui2025forada}, dual-stream adapter~\cite{liao2026dfa}, RepDFD~\cite{lin2025repdfd}, AdaptPrompt~\cite{jiang2026adaptprompt}, separable prompts~\cite{yang2026separable} & higher & Flexible; can inject forensic priors (boundaries, frequency, text). Adds inference cost and design choices. \\
Reparam. & Learn a structured update of frozen weights, e.g.\ low-rank $\Delta\mathbf{W}=\mathbf{B}\mathbf{A}$ or a residual singular-value subspace & LoRA~\cite{hu2022lora}$^\ast$, Effort (SVD residual)~\cite{yan2025effort}$^\ast$, environment-invariant subspaces~\cite{chen2026envsubspace} & same if merged & Mergeable after training; adaptation controlled by the rank. Larger budgets than selective tuning. \\
Hybrid & Combine a PEFT scheme with another mechanism: auxiliary objectives, mixtures of experts, temporal modules, weight averaging & GenD (LN + alignment/uniformity)~\cite{yermakov2025gend}$^\ast$, MoE of LoRA experts~\cite{liu2024mole}, shortcut suppression~\cite{wang2026shortcut}, dual supervision~\cite{zhang2026dino}, temporal adapter~\cite{han2024fcg}, soups / WiSE-FT~\cite{wortsman2022soups,wortsman2022wiseft}$^\ast$ & varies & Often the best reported numbers. More complex and harder to reproduce; the contribution of each part is rarely isolated. \\
\midrule
Bounds & Head only / all weights & Linear probe$^\ast$, full fine-tuning$^\ast$ & same & Under-adapts / forgets the pre-trained features. \\
\bottomrule
\end{tabular}
\end{table*}

\subsection{Selective methods}
Selective methods train a subset of the existing parameters and therefore add no inference cost. The extreme case trains only a classifier on frozen features: Ojha et al.~\cite{ojha2023universal} showed that a nearest-neighbour or linear classifier on frozen CLIP features generalizes across generators far better than detectors trained end to end; training-free variants score images by the stability of their features under perturbations~\cite{tsai2024trainingfree}, and patch-token probes use features below the class token~\cite{abdullah2026tap}. For face forgeries, where the manipulation is local and subtle, frozen features leave much of the attainable accuracy unused (Sec.~\ref{sec:bench}). Slightly larger subsets close this gap. Yermakov et al.~\cite{yermakov2025clip,yermakov2025gend} tune only the LayerNorm parameters of CLIP; BitFit~\cite{benzaken2022bitfit} tunes all biases, and BPL~\cite{xu2026bpl} tunes only biases while training on mined real--fake pairs; BiasLN tunes both biases and LayerNorm scales. Zhang et al.~\cite{zhang2026layermask} select which layers to adapt according to forgery-relevant activations. Bias and LayerNorm parameters change only per-channel offsets and scales of the frozen computation; Sec.~\ref{sec:analysis} makes precise how they overlap and why they differ.

\subsection{Additive methods}
Additive methods keep the encoder intact and add capacity next to it. Forensics Adapter~\cite{cui2025forada} attaches a lightweight adapter network to CLIP that learns blending boundaries and interacts with the frozen encoder; Liao et al.~\cite{liao2026dfa} use a dual-stream variant. Prompt-based methods learn input tokens instead: RepDFD~\cite{lin2025repdfd} reprograms CLIP through input perturbations and face-aware text prompts, AdaptPrompt~\cite{jiang2026adaptprompt} learns visual and textual prompts, and Yang et al.~\cite{yang2026separable} exploit the text encoder with separable prompts. The added modules make it easy to inject forensic priors, such as boundary or frequency cues, at the price of extra inference computation and more design choices.

\subsection{Reparameterization methods}
Reparameterization methods learn a structured update of the frozen weights. LoRA~\cite{hu2022lora} adds a trainable low-rank product $\mathbf{B}\mathbf{A}$ to selected weight matrices. Effort~\cite{yan2025effort} decomposes each attention projection by singular-value decomposition, freezes the principal components that carry the pre-trained semantics and tunes a small residual subspace, with an orthogonality constraint. Related methods separate forgery-specific from environment-specific directions~\cite{chen2026envsubspace}. All of these updates can be merged into the weights after training; implementations that rebuild the weights at every forward pass add latency (Sec.~\ref{sec:cost}).

\subsection{Hybrid methods}
Hybrid methods pair a PEFT scheme with a second mechanism. GenD~\cite{yermakov2025gend} tunes the LayerNorm parameters, normalizes the features onto the hypersphere and adds alignment and uniformity losses~\cite{wang2020alignment}. Liu et al.~\cite{liu2024mole} combine several LoRA experts in a mixture for transferable AI-generated image detection. Other works add objectives that suppress forgery-specific shortcuts~\cite{wang2026shortcut} or supervise both real/fake labels and finer forgery categories on a DINO backbone~\cite{zhang2026dino}. For video, Han et al.~\cite{han2024fcg} insert a temporal adapter guided by facial components. Post-hoc averaging of fine-tuned weights (model soups~\cite{wortsman2022soups}) or interpolation towards the pre-trained weights (WiSE-FT~\cite{wortsman2022wiseft}) improves robustness in image classification without extra inference cost; its effect on PEFT detectors had not been studied. Hybrids often report the best numbers, but the share of each component is rarely measured; our benchmark isolates it for GenD and for weight averaging.

\subsection{Beyond the encoder adaptation}
Data-centric methods synthesize training forgeries such as self-blended images~\cite{shiohara2022sbi,chen2022sladd,lin2024cdfa} and are complementary to all branches above. Design choices of the training pipeline (number of epochs, augmentation) can matter as much as the architecture~\cite{pellegrini2025design}, intermediate encoder layers can generalize better than the last one~\cite{park2025mold}, and recent studies report clear cross-domain limits of foundation-model detectors~\cite{delibasoglu2026limits}.

% =====================================================================
\section{How Methods Are Evaluated}
\label{sec:eval}

\subsection{Checkpoint selection}
For a detector trained on FF++, the epoch whose checkpoint is reported strongly affects cross-dataset numbers, because in-domain and out-of-domain accuracy do not peak together. A clean protocol selects the checkpoint on a validation split of the \emph{training} distribution (here, FF++), so that the cross-dataset benchmarks remain unseen. Several widely used protocols instead select on data from the \emph{test} benchmarks:
\begin{itemize}[leftmargin=*, topsep=2pt, itemsep=2pt, label=--]
    \item The DeepfakeBench trainer~\cite{yan2023deepfakebench}, on which several detectors, including Effort~\cite{yan2025effort}, build, evaluates every test dataset after each epoch and keeps a best checkpoint \emph{per test dataset}; the reported number of each test set then comes from the epoch that is best on that very test set.
    \item GenD~\cite{yermakov2025gend} selects on a custom validation set drawn from several datasets, some of which are also its test benchmarks (with disjoint samples), so the validation data share the target domains.
    \item Many papers do not state the rule.
\end{itemize}
Selection on test-benchmark data turns part of the test benchmark into a validation set and makes cross-dataset results optimistic (Fig.~\ref{fig:protocol}). Sec.~\ref{sec:selection} quantifies the effect for one model.

\begin{figure*}[!t]
\centering
\includegraphics[width=0.9\textwidth]{figs/fig_protocol_figma.pdf}
\caption{Clean checkpoint selection on FF++ validation data (green, our benchmark) versus selection on test-benchmark data (red), which lets test data influence the reported model.}
\label{fig:protocol}
\end{figure*}

\subsection{Averages, frames and seeds}
Cross-dataset averages are computed over different dataset sets (with or without UADFV, for example), so averages from different papers are not comparable even when the per-dataset numbers are. Frame- and video-level AUC are mixed, and the number of training frames per video varies from 8 to all 32 extracted frames. Finally, most methods are reported from a single training run, while our runs show seed-to-seed standard deviations of about 0.005 AUC for the cross-dataset average, larger than many reported improvements. Table~\ref{tab:literature} in the Appendix lists published numbers for reference; they are copied from the respective papers (mostly from Table~1 of Effort~\cite{yan2025effort}), were obtained under the different rules described above and are \emph{not} comparable with our controlled runs.



% =====================================================================
\section{Controlled Benchmark}
\label{sec:bench}

\subsection{Protocol}
\label{sec:protocol}
All configurations use CLIP ViT-L/14 with an $L_2$-normalized class token and a linear head, trained on FF++ c23 (official 720/140/140 split, 8 evenly spaced frames per video unless stated) with the DeepfakeBench pre-processing and augmentation (flip, rotation, blur, colour jitter, JPEG). Training runs for 10 epochs with Adam, cosine learning-rate decay and batch size 32 in bfloat16 on one A100. \textbf{Every checkpoint is selected by frame AUC on the FF++ validation split}, which shares no video or frame with any test set. Testing uses all 32 frames per video and reports video AUC; Avg(D) is the mean over CDF-v2, DFD, DFDC and DFDCP, and Avg(M) over the seven DF40 face-swapping sets (UniFace, BlendFace, E4S, FaceDancer, FSGAN, InSwap, SimSwap) applied to FF++ identities. An automated audit confirms that the validation split shares no frame and no video with any test set. The loss is the class-weighted cross-entropy with label smoothing that the code implements,
\begin{equation}
\mathcal{L}=-\frac{1}{\sum_i w_{y_i}}\sum_{i} w_{y_i}\sum_{c\in\{0,1\}}\tilde{y}_{i,c}\log p_{i,c},
\label{eq:ce}
\end{equation}
where $\tilde{y}_{i,c}=(1-\alpha)\,\mathbf{1}[y_i=c]+\alpha/2$, with labels $0$ (real) and $1$ (fake), weights $w_0=1$, $w_1=2$ and $\alpha=0.1$; GenD adds its own losses to Eq.~\eqref{eq:ce}, and no other scheme uses an auxiliary loss. Bias tuning, LayerNorm tuning and BiasLN are run with ten seeds (1024 and 1--9); GenD and BiasLN with a larger LayerNorm-scale learning rate with six (1024 and 1--5); the other schemes with seeds 1024, 1 and 2. Pairwise tests use the seeds that both schemes share. We report the mean and sample standard deviation and compare methods with a paired $t$-test over seeds and a two one-sided tests (TOST) equivalence test with a margin of $\pm0.01$ AUC. For the seven main comparisons on the two averages (14 tests), we correct both the $t$-test and the TOST $p$-values with the Holm procedure (Table~\ref{tab:stats}); single-dataset tests in the text are uncorrected and are marked as such.

Re-trained schemes: linear probe; full fine-tuning (encoder learning rate $10^{-5}$); BitFit (all biases, including LayerNorm shifts); LayerNorm tuning; BiasLN (biases and LayerNorm scales); LoRA ($r{=}8$ on query and value); Effort (rank-1 SVD residual on all attention projections, following its released code); and GenD (LayerNorm tuning with alignment and uniformity losses, following its released code).

\subsection{Effect of the selection rule}
\label{sec:selection}
The same BiasLN configuration reaches an Avg(D) of 0.934 when its checkpoint is selected on a subset of CDF-v2 and \num{AvgDBiasln} when selected on FF++ validation. This difference of about one point exceeds the gap between all PEFT schemes in Table~\ref{tab:main}. Selection on FF++ validation is itself imperfect: its AUC saturates above 0.98 and does not rank checkpoints by their cross-dataset accuracy, so a validation set built from non-test sources is an open need (Sec.~\ref{sec:open}).

\begin{table*}[!t]
\centering
\caption{Video AUC of re-trained schemes under one protocol (FF++ c23 training, checkpoint selected on FF++ validation). Avg(D) and Avg(M) are mean $\pm$ std over $n$ seeds; per-dataset values are seed means. Bold: best; shaded: second best.}
\label{tab:main}
\resizebox{\textwidth}{!}{\input{generated/tab_main.tex}}
\vspace{3mm}
\caption{Ablation of selective and low-rank schemes and of $L_2$ normalization before the classifier (same protocol). Bold: best; shaded: second best.}
\label{tab:ablation}
\resizebox{0.85\textwidth}{!}{\input{generated/tab_ablation.tex}}
\end{table*}

\subsection{RQ1: Adaptation versus preservation}
\label{sec:main}
Table~\ref{tab:main} and Table~\ref{tab:ablation} give the main results; Tables~\ref{tab:perdataD} and~\ref{tab:perdataM} repeat them per dataset with standard deviations, and Fig.~\ref{fig:tradeoff} plots accuracy against the number of trainable parameters.
\textbf{Adaptation is essential but light.} The linear probe, which preserves the encoder completely, reaches only \num{AvgDLinprobe} Avg(D); full fine-tuning, which adapts everything, reaches \num{AvgDFullft} although its in-domain AUC is high (\num{AvgFFFullft} on FF++). Every PEFT scheme lies between 0.90 and 0.93. The trade-off is therefore not monotonic: some adaptation is needed to expose local forgery cues, but beyond about 0.1--0.3M parameters additional capacity is spent on fitting the training forgeries.
\textbf{More parameters do not help.} LoRA trains $2.4\times$ more parameters than BiasLN and $7.6\times$ more than LayerNorm tuning and is not better on either average, with the largest seed variance of all PEFT schemes; Effort is about one point below the affine schemes on Avg(D) ($p=0.10$ against BiasLN and $p=0.003$ against GenD, three seeds, uncorrected) and comparable on Avg(M). In Fig.~\ref{fig:tradeoff}, the highest Avg(M) is reached at the smallest selective budget (0.1M), and no larger budget exceeds it on Avg(D) beyond the seed variation.
\textbf{On the cross-dataset average the affine schemes are equivalent.} Bias tuning (\num{AvgDBitfit}), LayerNorm tuning (\num{AvgDLnLC}), BiasLN (\num{AvgDBiasln}) and GenD (\num{AvgDGend}) differ by less than the seed-to-seed variation; none of the pairwise $t$-tests is significant, and after Holm correction every pair of the four schemes is equivalent within $\pm0.01$ AUC ($p^{\mathrm{Holm}}_{\mathrm{TOST}}\le0.047$, Table~\ref{tab:stats}). Differences appear per dataset, which RQ2 and RQ3 examine.

\begin{table*}[!t]
\centering
\caption{Paired comparisons over the seeds shared by both schemes. $\Delta$: mean difference A$-$B; $p_{\mathrm{Holm}}$: Holm-corrected over the 14 $t$-tests; $p^{\mathrm{Holm}}_{\mathrm{TOST}}$: Holm-corrected equivalence test ($\pm0.01$ AUC; small means equivalent); wins: seeds where A is better. Bold: $<0.05$.}
\label{tab:stats}
\resizebox{\textwidth}{!}{\input{generated/tab_stats.tex}}
\end{table*}

\begin{figure}[!t]
\centering
\includegraphics[width=\columnwidth]{generated/fig_params_auc.pdf}
\caption{Avg(D) (top) and Avg(M) (bottom) against trainable parameters (log scale), mean $\pm$ std over seeds; insets enlarge the PEFT region.}
\label{fig:tradeoff}
\end{figure}

\begin{table}[!t]
\centering
\caption{Cross-dataset video AUC ($\times100$), mean $\pm$ std over $n$ seeds (same runs as Table~\ref{tab:main}). Bold: best; shaded: second best.}
\label{tab:perdataD}
\resizebox{\columnwidth}{!}{\input{generated/tab_perdata_D.tex}}
\end{table}

\begin{table*}[!t]
\centering
\caption{Cross-method video AUC ($\times100$) on the seven DF40 face-swapping sets, mean $\pm$ std over $n$ seeds. Bold: best; shaded: second best.}
\label{tab:perdataM}
\resizebox{\textwidth}{!}{\input{generated/tab_perdata_M.tex}}
\end{table*}

\begin{figure}[!t]
\centering
\includegraphics[width=\columnwidth]{generated/fig_perdataset.pdf}
\caption{Cross-dataset video AUC (mean $\pm$ std over seeds).}
\label{fig:perdata}
\end{figure}



\subsection{RQ2: Where and what to adapt}
\label{sec:ablation}
\textbf{Biases versus LayerNorm parameters.} The two groups of per-channel parameters behave differently per dataset (Tables~\ref{tab:perdataD} and~\ref{tab:perdataM}, Fig.~\ref{fig:perdata}). Adding the biases to LayerNorm tuning (BiasLN) raises in-domain AUC on FF++ from \num{AvgFFLnLC} to \num{AvgFFBiasln} ($p<0.001$; per manipulation in Table~\ref{tab:indomain}, Appendix) and DFDCP by 3.2 points (both in ten of ten seeds, $p<0.001$ uncorrected), but lowers DFDC by 1.7 points (ten of ten seeds, $p<0.001$ uncorrected) and Avg(M) by 0.006 (nine of ten seeds, $p=0.003$, $p_{\mathrm{Holm}}=0.048$). The biases thus trade transfer to new manipulation methods for accuracy on data closer to the training distribution, and BiasLN is equivalent to bias tuning alone on both averages ($p^{\mathrm{Holm}}_{\mathrm{TOST}}<0.001$). Sec.~\ref{sec:analysis} relates this to the structure of the two parameter groups.

\textbf{Roles and depth.} Table~\ref{tab:roles} splits the tuned set by role, and Table~\ref{tab:depth} by depth. No single role is clearly better: query and fc1 biases, the residual-stream biases (value, output, fc2), LayerNorm scales or shifts alone all reach 0.914--0.920 Avg(D). Depth matters more. Tuning only six or twelve blocks loses one to three points, wherever the blocks are; tuning blocks 0--18 matches tuning all 24 blocks with 20\% fewer parameters. Restricting only the biases to blocks 0--18 while tuning the LayerNorm parameters everywhere does not recover the cross-method transfer of LayerNorm tuning (Avg(M) 0.957 over two seeds), so the effect of the biases is not confined to the last blocks. A gradient profile confirms that the key biases receive exactly zero gradient (Proposition~\ref{prop:key}) and that the last blocks receive the smallest updates.

\begin{table}[!t]
\centering
\caption{Tuned subset inside BiasLN ($L_2$ normalization, cross-entropy). Video AUC, mean $\pm$ std over $n$ seeds.}
\label{tab:roles}
\resizebox{\columnwidth}{!}{\input{generated/tab_roles.tex}}
\end{table}

\begin{table}[!t]
\centering
\caption{Depth of adaptation: BiasLN restricted to a range of transformer blocks (24 in total). Seed 1024 unless $n$ is given.}
\label{tab:depth}
\footnotesize
\begin{tabular}{l r c c c}
\toprule
\textbf{Blocks} & \textbf{Params} & $n$ & \textbf{Avg(D)} & \textbf{Avg(M)} \\
\midrule
0--5 & 86k & 1 & 0.894 & 0.939 \\
18--23 & 86k & 1 & 0.885 & 0.945 \\
0--11 & 166k & 1 & 0.904 & 0.946 \\
12--23 & 166k & 1 & 0.891 & 0.939 \\
12--18 & 99k & 1 & 0.903 & 0.945 \\
6--17 & 166k & 3 & 0.919\,{\scriptsize$\pm$0.002} & 0.966\,{\scriptsize$\pm$0.004} \\
0--18 & 259k & 3 & 0.923\,{\scriptsize$\pm$0.004} & 0.960\,{\scriptsize$\pm$0.005} \\
0--23 (all) & 326k & \num{NDBiasln} & \num{AvgDBiasln}\,{\scriptsize$\pm$\num{StdDBiasln}} & \num{AvgMBiasln}\,{\scriptsize$\pm$\num{StdMBiasln}} \\
\bottomrule
\end{tabular}
\end{table}

\subsection{RQ3: Unseen generators}
\label{sec:rq3}
The seven DF40 sets apply newer face-swapping generators (UniFace, BlendFace, E4S, FaceDancer, FSGAN, InSwap, SimSwap) to the FF++ identities, so they isolate the change of generator from the change of identities and capture conditions. Two observations stand out (Table~\ref{tab:perdataM}). First, transfer to new generators is not ordered by budget: tuning only the LayerNorm parameters (0.1M) reaches \num{AvgMLnLC} Avg(M), above BiasLN ($+0.006$, nine of ten seeds, $p_{\mathrm{Holm}}=0.048$) and bias tuning ($+0.006$, eight of ten seeds; $p=0.020$ uncorrected, $p_{\mathrm{Holm}}=0.26$), with the largest gains on SimSwap ($+2.1$ points against BiasLN in every seed), UniFace and InSwap, while LoRA (0.8M) is lowest among the PEFT schemes. Second, the auxiliary losses of the best hybrid scheme add little: LayerNorm tuning tunes the same parameters as GenD without its alignment and uniformity losses and is equivalent to it on Avg(M) ($-0.001$, $p^{\mathrm{Holm}}_{\mathrm{TOST}}=0.016$) and on Avg(D) ($-0.003$, $0.047$). The losses help mainly on DFDCP ($+0.010$, $p=0.019$ uncorrected) and lower the seed variance. BlendFace remains the hardest set for every scheme.

Our benchmark does not cover entire-face synthesis or diffusion-generated images. In the literature, frozen CLIP features transfer from GAN to diffusion images with a linear probe~\cite{ojha2023universal}, and mixtures of LoRA experts have been proposed for transferable AI-generated image detection~\cite{liu2024mole}, but PEFT schemes have not been compared on such data under a common protocol (Sec.~\ref{sec:open}).

\subsection{Optimization, data and weight averaging}
\label{sec:extra}
Table~\ref{tab:extra} (Appendix) collects variants that change the training or post-processing rather than the tuned set. Raising the learning rate of the LayerNorm scales of BiasLN by $3\times$ changes Avg(M) by $+0.003$ over six seeds, which is not significant ($p=0.29$), and Avg(D) by $-0.002$; a $10\times$ rate hurts. How fast the scales move is therefore not what separates BiasLN from LayerNorm tuning. Training on all 32 frames per video (115k images instead of 29k) helps neither BiasLN nor GenD. Concatenating class tokens from blocks 12, 18 and 24 does not help. Uniform averaging of the last epochs or of three seeds changes the averages by less than the seed variation, and WiSE-FT interpolation with the pre-trained weights collapses (Avg(D) 0.82--0.86) because the linear head no longer matches the interpolated features.



\subsection{Computational cost}
\label{sec:cost}
Table~\ref{tab:cost} reports measured costs. Training costs of all PEFT schemes are within 15\% of each other, because back-propagation through the frozen 24-block encoder dominates; full fine-tuning needs 62\% more memory. The schemes differ at inference: selective tuning keeps the cost of the frozen encoder (162 GFLOPs per image), while Effort and unmerged LoRA rebuild or add weights in every forward pass.

\begin{table}[!t]
\centering
\caption{Measured cost on an NVIDIA A100-40GB (fp32, $224\times224$). Train: batch 32, forward + backward + Adam step. Infer: batch 1. Ckpt: size of the stored trainable tensors.}
\label{tab:cost}
\resizebox{\columnwidth}{!}{%
\begin{tabular}{l r r r r r r}
\toprule
\textbf{Method} & \textbf{Trainable} & \textbf{Ckpt} & \textbf{Train} & \textbf{Train} & \textbf{Infer} & \textbf{GFLOPs} \\
 & & \textbf{(MB)} & \textbf{img/s} & \textbf{mem (GB)} & \textbf{ms} & \\
\midrule
Linear probe & 0.002M & 0.0 & 100.6 & 1.67 & 17.0 & 162.0 \\
LayerNorm (GenD) & 0.104M & 0.4 & 47.1 & 12.16 & 17.1 & 162.0 \\
Effort & 0.199M & 0.8 & 41.7 & 13.22 & 23.8 & 162.2 \\
BitFit & 0.274M & 1.1 & 46.6 & 12.16 & 17.1 & 162.0 \\
BiasLN & 0.326M & 1.3 & 46.5 & 12.16 & 17.2 & 162.0 \\
LoRA $r{=}8$ & 0.788M & 3.0 & 45.9 & 12.87 & 18.6 & 162.4 \\
Full FT & 303.2M & 1156.7 & 32.8 & 19.72 & 17.2 & 162.0 \\
\bottomrule
\end{tabular}}
\end{table}

% =====================================================================
\section{What Bias and LayerNorm Tuning Can Express}
\label{sec:analysis}

Bias tuning, LayerNorm tuning and their union BiasLN act on per-channel quantities of the frozen computation. Two facts relate them.

\begin{proposition}[Key bias is redundant~\cite{namazifar2023bias}]
\label{prop:key}
In softmax attention with scores $s_{ij}=\mathbf{q}_i^\top(\mathbf{W}_k\mathbf{u}_j+\mathbf{b}_k)/\sqrt{d}$, the term $\mathbf{q}_i^\top\mathbf{b}_k/\sqrt{d}$ is the same for every key $j$ and cancels in the softmax; hence $\partial\mathcal{L}/\partial\mathbf{b}_k=\mathbf{0}$.
\end{proposition}

\begin{proposition}[LayerNorm shift is a tied bias update]
\label{prop:fold}
In a pre-norm block, changing the shift of a LayerNorm by $\Delta\boldsymbol{\beta}$ changes the following linear layer exactly as changing its bias by $\mathbf{W}\Delta\boldsymbol{\beta}$, since $\mathbf{W}(\boldsymbol{\gamma}\odot\hat{\mathbf{u}}+\boldsymbol{\beta}+\Delta\boldsymbol{\beta})+\mathbf{b}=\mathbf{W}(\boldsymbol{\gamma}\odot\hat{\mathbf{u}}+\boldsymbol{\beta})+(\mathbf{b}+\mathbf{W}\Delta\boldsymbol{\beta})$.
\end{proposition}

Consequently, LayerNorm-shift tuning is contained in bias tuning, and the LayerNorm scales $\boldsymbol{\gamma}$ are the only per-channel multiplicative parameters; BiasLN is the complete per-channel affine adaptation of the encoder. The containment is strict in a useful sense. The first LayerNorm of a block feeds the query, key and value projections through one shared $\Delta\boldsymbol{\beta}\in\mathbb{R}^{d}$, so the induced bias updates $(\mathbf{W}_q\Delta\boldsymbol{\beta},\mathbf{W}_k\Delta\boldsymbol{\beta},\mathbf{W}_v\Delta\boldsymbol{\beta})$ are tied to one another and to the pre-trained weights, whereas bias tuning moves the $3d$ attention biases and the $4d$ MLP bias independently. Likewise, the second LayerNorm reaches the MLP only through $\mathbf{W}_{\mathrm{fc1}}$, while bias tuning also shifts the residual stream directly through $\mathbf{b}_o$ and $\mathbf{b}_{\mathrm{fc2}}$.

LayerNorm tuning is therefore a lower-dimensional adaptation whose updates follow directions that the pre-trained encoder already uses, and the results are consistent with this view. With the free biases, the model fits the training distribution better: it has the highest in-domain AUC among the affine schemes (Table~\ref{tab:indomain}) and the highest DFDCP score, a dataset close to FF++ in capture conditions. Without them, the adaptation tends to transfer better to new face-swapping generators applied to the same identities (significant against BiasLN after multiple-comparison correction, consistent against bias tuning) (Avg(M), Table~\ref{tab:perdataM}). Because the scales are already present in LayerNorm tuning, adding the biases to it (BiasLN) moves the model towards the bias-tuning behaviour, and BiasLN is statistically equivalent to bias tuning. On the cross-dataset average the two effects cancel, which is why all affine schemes look alike in Avg(D) alone. This is a consistency argument rather than a proof; testing it on other backbones and training sets is left open (Sec.~\ref{sec:open}).

% =====================================================================
\section{Recommendations}
\label{sec:recommend}
\begin{enumerate}[leftmargin=*, topsep=2pt, itemsep=2pt]
    \item \textbf{State and fix the selection rule}, and never select on data from a test benchmark. Report the selected epoch.
    \item \textbf{Report several seeds} with standard deviations, and test differences with paired or equivalence tests; differences below about 0.01 Avg(D) are within seed noise for CLIP ViT-L/14 on FF++.
    \item \textbf{Re-train baselines} under the same pipeline instead of copying published numbers, and label copied numbers as not comparable.
    \item \textbf{Fix the test-set average} (we follow Effort~\cite{yan2025effort}: CDF-v2, DFD, DFDC, DFDCP) and report per-dataset video AUC.
    \item \textbf{Report measured costs}, including inference latency; parameter counts alone hide the cost of unmerged reparameterizations and added modules.
    \item As a strong and simple baseline, \textbf{tune the LayerNorm parameters} of the encoder with an $L_2$-normalized linear head: it is equivalent to GenD on both averages without auxiliary losses. Add the biases (BiasLN) only when in-domain accuracy or data close to the training distribution matters more than transfer to new manipulation methods, and \textbf{report per-dataset results}, since the averages hide this trade-off.
\end{enumerate}

% =====================================================================
\section{Future Directions}
\label{sec:open}
\textbf{Continual learning with PEFT.} New generators appear faster than detectors can be re-trained. Because PEFT modules are small and separable, a natural direction is to add one module (a LoRA, an adapter or a set of LayerNorm parameters) per new forgery family while keeping the shared encoder and the earlier modules fixed, so that detection of older forgeries is not forgotten; how to route between such modules at test time, and whether selective schemes such as LayerNorm tuning are less prone to forgetting than additive ones, is open.
\textbf{PEFT for video and multimodal detection.} Most PEFT detectors score frames independently. Temporal adapters~\cite{han2024fcg} add dedicated modules, while simple temporal pooling of PEFT features and audio--visual PEFT, which would exploit lip--speech inconsistencies, remain largely unexplored.
\textbf{Unseen generator families.} Our benchmark and most PEFT detectors are trained and tested on face swapping and reenactment. A controlled comparison of PEFT branches from face-swap training to diffusion-generated and entire-face synthesis, and in the opposite direction, is needed to answer RQ3 beyond face swapping.
\textbf{Validation without test leakage.} FF++ validation saturates; a validation set built from manipulations that are not used for testing (for example, the reenactment subsets of DF40) could select better checkpoints without using test benchmarks.
\textbf{Scope of the bias--LayerNorm trade-off and other backbones.} Our evidence comes from one backbone (CLIP ViT-L/14) and one training set (FF++); whether the advantage of LayerNorm tuning on unseen generators holds, with more seeds, on other encoders such as DINO~\cite{caron2021emerging}, and with synthetic training forgeries~\cite{shiohara2022sbi}, is untested.
\textbf{Deployment.} A ViT-L/14 encoder needs 162 GFLOPs per image; distilling PEFT detectors into smaller encoders and testing robustness to real-world post-processing are needed for practical use.

\paragraph{Limitations of this study.}\label{sec:limits} The benchmark scores individual frames of pre-cropped faces and averages them per video, without temporal modelling. It covers face swapping and reenactment in FF++, four cross-dataset benchmarks and seven DF40 face-swapping methods, but not entire-face synthesis or diffusion-based editing. We measured latency, memory and FLOPs on a data-centre GPU, not energy or edge hardware. Six to ten seeds give sufficient power for the average AUCs, but per-dataset differences below about 0.005 remain within seed noise.

% =====================================================================
\section{Conclusion}
We surveyed parameter-efficient fine-tuning for face deepfake detection, organized existing detectors into selective, additive, reparameterization and hybrid branches, and reviewed how they are evaluated, finding that several common protocols select checkpoints on data from the test benchmarks. A controlled benchmark under a leakage-free protocol answers our three research questions. Light adaptation is essential but sufficient, and larger budgets do not help (RQ1). Which parameters are adapted matters more than how many: adding LayerNorm scales to bias tuning changes nothing measurable, and tuning only the LayerNorm parameters transfers better to unseen face-swapping methods than BiasLN, whereas the biases raise in-domain accuracy (RQ2). Transfer to new generators is not ordered by budget, and the auxiliary losses of the strongest hybrid scheme add little once the same LayerNorm parameters are tuned (RQ3). We hope that the taxonomy, the released protocol and the per-run results help future work compare PEFT detectors on equal terms.

\appendix
\renewcommand{\thesection}{A}\renewcommand{\thetable}{A\arabic{table}}\setcounter{table}{0}
\section{Published Results}
\label{sec:appendix}
Table~\ref{tab:literature} collects the video AUC reported by the authors of detectors trained on FF++ c23. It is given for orientation only; Sec.~\ref{sec:eval} explains why these numbers cannot be compared with each other or with our controlled runs.

\begin{table*}[!t]
\centering
\caption{Published video AUC of detectors trained on FF++ c23, as reported by their authors (mostly Table~1 of Effort~\cite{yan2025effort}). Selection rules, frames and seeds differ between rows; the numbers are listed for reference only and are not comparable with Tables~\ref{tab:main}--\ref{tab:perdataM}. $^\dagger$: values reported by Effort.}
\label{tab:literature}
\resizebox{\textwidth}{!}{\input{generated/tab_literature.tex}}
\end{table*}

Table~\ref{tab:indomain} gives the in-domain AUC on the four FF++ manipulations (Sec.~\ref{sec:ablation}), and Table~\ref{tab:extra} the training and post-hoc variants of Sec.~\ref{sec:extra}.

\begin{table}[!t]
\centering
\caption{In-domain video AUC ($\times100$) on the four FF++ manipulations (test split), seed means.}
\label{tab:indomain}
\resizebox{\columnwidth}{!}{\input{generated/tab_indomain.tex}}
\end{table}

\begin{table}[!t]
\centering
\caption{Training and post-hoc variants (video AUC). Seed 1024 unless $n$ is given. Soups average seeds 1024, 1 and 2.}
\label{tab:extra}
\footnotesize
\resizebox{\columnwidth}{!}{%
\begin{tabular}{l c c c}
\toprule
\textbf{Variant} & $n$ & \textbf{Avg(D)} & \textbf{Avg(M)} \\
\midrule
BiasLN (reference) & \num{NDBiasln} & \num{AvgDBiasln}\,{\scriptsize$\pm$\num{StdDBiasln}} & \num{AvgMBiasln}\,{\scriptsize$\pm$\num{StdMBiasln}} \\
LN-scale lr $\times3$ & 6 & \num{AvgDGlrD}\,{\scriptsize$\pm$\num{StdDGlrD}} & \num{AvgMGlrD}\,{\scriptsize$\pm$\num{StdMGlrD}} \\
LN-scale lr $\times10$ & 1 & 0.904 & 0.961 \\
Multi-layer CLS (12+18+24) & 1 & 0.917 & 0.959 \\
Blocks 0--18, 8 frames & 1 & 0.924 & 0.956 \\
Blocks 0--18, 32 frames & 1 & 0.917 & 0.959 \\
GenD, 8 frames & 1 & 0.925 & 0.966 \\
GenD, 32 frames & 1 & 0.924 & 0.963 \\
Seed soup, BiasLN & -- & 0.920 & 0.939 \\
Seed soup, GenD & -- & 0.927 & 0.966 \\
Epoch soup ($\geq5$), BiasLN & 2 & 0.921\,{\scriptsize$\pm$0.004} & 0.962\,{\scriptsize$\pm$0.004} \\
WiSE-FT $\alpha{=}0.5$, BiasLN & -- & 0.826 & 0.829 \\
\bottomrule
\end{tabular}}
\end{table}

\begin{thebibliography}{99}
\balance

\bibitem{goodfellow2020gan}
Goodfellow I, Pouget-Abadie J, Mirza M, Xu B, Warde-Farley D, Ozair S, Courville A, Bengio Y.
Generative adversarial networks.
Commun ACM. 2020;63(11):139--44.

\bibitem{yang2022diffusion}
Yang L, Zhang Z, Song Y, Hong S, Xu R, Zhao Y, Zhang W, Cui B, Yang MH.
Diffusion models: A comprehensive survey of methods and applications.
ACM Comput Surv. 2023;56(4):1--39.

\bibitem{pei2024survey}
Pei G, Zhang J, Hu M, Zhang Z, Wang C, Wu Y, Zhai G, Yang J, Tao D.
Deepfake generation and detection: A benchmark and survey.
ACM Comput Surv. arXiv:2403.17881. 2024.

\bibitem{liu2024evolving}
Liu P, Tao Q, Zhou JT.
Evolving from single-modal to multi-modal facial deepfake detection: Progress and challenges.
arXiv preprint arXiv:2406.06965. 2024.

\bibitem{rossler2019ff}
R{\"o}ssler A, Cozzolino D, Verdoliva L, Riess C, Thies J, Nie{\ss}ner M.
FaceForensics++: Learning to detect manipulated facial images.
In: Proc IEEE/CVF Int Conf Comput Vis (ICCV). 2019.

\bibitem{qian2020frequency}
Qian Y, Yin G, Sheng L, Chen Z, Shao J.
Thinking in frequency: Face forgery detection by mining frequency-aware clues.
In: Proc Eur Conf Comput Vis (ECCV). 2020.

\bibitem{luo2021srm}
Luo Y, Zhang Y, Yan J, Liu W.
Generalizing face forgery detection with high-frequency features.
In: Proc IEEE/CVF Conf Comput Vis Pattern Recognit (CVPR). 2021.

\bibitem{yan2023deepfakebench}
Yan Z, Zhang Y, Yuan X, Lyu S, Wu B.
DeepfakeBench: A comprehensive benchmark of deepfake detection.
Adv Neural Inf Process Syst (NeurIPS), Datasets and Benchmarks Track. 2023;36:4534--65.

\bibitem{radford2021clip_vitl}
Radford A, Kim JW, Hallacy C, Ramesh A, Goh G, Agarwal S, et al.
Learning transferable visual models from natural language supervision.
In: Proc Int Conf Mach Learn (ICML). 2021;8748--63.

\bibitem{ojha2023universal}
Ojha U, Li Y, Lee YJ.
Towards universal fake image detectors that generalize across generative models.
In: Proc IEEE/CVF Conf Comput Vis Pattern Recognit (CVPR). 2023;24480--9.

\bibitem{tsai2024trainingfree}
Tsai CT, Ko CY, Chung IH, Wang YCF, Chen PY.
Understanding and improving training-free AI-generated image detections with vision foundation models.
arXiv preprint arXiv:2411.19117. 2024.

\bibitem{yan2025effort}
Yan Z, Wang J, Jin P, Zhang KY, Liu C, Chen S, Yao T, Ding S, Wu B, Yuan L.
Orthogonal subspace decomposition for generalizable AI-generated image detection.
In: Proc Int Conf Mach Learn (ICML). 2025.

\bibitem{yermakov2025gend}
Yermakov A, Cech J, Matas J, Fritz M.
Deepfake detection that generalizes across benchmarks.
In: Proc IEEE/CVF Winter Conf Appl Comput Vis (WACV). 2026;773--83.

\bibitem{yermakov2025clip}
Yermakov A, Cech J, Matas J.
Unlocking the hidden potential of CLIP in generalizable deepfake detection.
arXiv preprint arXiv:2503.19683. 2025.

\bibitem{cui2025forada}
Cui X, Li Y, Zhu D, Zhou J, Dong J, Lyu S.
Forensics Adapter: Adapting CLIP for generalizable face forgery detection.
In: Proc IEEE/CVF Conf Comput Vis Pattern Recognit (CVPR). 2025;19207--17.

\bibitem{kumar2022finetuning}
Kumar A, Raghunathan A, Jones R, Ma T, Liang P.
Fine-tuning can distort pretrained features and underperform out-of-distribution.
In: Int Conf Learn Represent (ICLR). 2022.

\bibitem{li2020celeb}
Li Y, Yang X, Sun P, Qi H, Lyu S.
Celeb-DF: A large-scale challenging dataset for deepfake forensics.
In: Proc IEEE/CVF Conf Comput Vis Pattern Recognit (CVPR). 2020;3207--16.

\bibitem{dufour2019dfd}
Dufour N, Gully A.
Contributing data to deepfake detection research.
Google Research Blog. 2019.

\bibitem{dolhansky2019dfdcp}
Dolhansky B, Howes R, Pflaum B, Baram N, Ferrer CC.
The deepfake detection challenge (DFDC) preview dataset.
arXiv preprint arXiv:1910.08854. 2019.

\bibitem{dolhansky2020dfdc}
Dolhansky B, Bitton J, Pflaum B, Lu J, Howes R, Wang M, Ferrer CC.
The DeepFake Detection Challenge (DFDC) dataset.
arXiv preprint arXiv:2006.07397. 2020.

\bibitem{yan2024df40}
Yan Z, Yao T, Chen S, Zhao Y, Fu X, Zhu J, Luo D, Wang C, Ding S, Wu Y, Yuan L.
DF40: Toward next-generation deepfake detection.
Adv Neural Inf Process Syst (NeurIPS), Datasets and Benchmarks Track. 2024.

\bibitem{caron2021emerging}
Caron M, Touvron H, Misra I, J{\'e}gou H, Mairal J, Bojanowski P, Joulin A.
Emerging properties in self-supervised vision transformers.
In: Proc IEEE/CVF Int Conf Comput Vis (ICCV). 2021.

\bibitem{zhang2026dino}
Zhang T, Yu P, Xia Z, Dai L, Zhou X, Gao H.
Fine-grained DINO tuning with dual supervision for face forgery detection.
In: Proc AAAI Conf Artif Intell (AAAI). 2026. arXiv:2511.12107.

\bibitem{zou2026clipdino}
Zou X, Zhou Y, Li K, Zhang S, Chen Y, Lang C, Xing J.
Harnessing CLIP and DINO: An uncertainty-aware cascaded fusion network for generalizable deepfake image detection.
arXiv preprint arXiv:2609.07670. 2026.

\bibitem{benzaken2022bitfit}
Ben Zaken E, Goldberg Y, Ravfogel S.
BitFit: Simple parameter-efficient fine-tuning for transformer-based masked language models.
In: Proc 60th Annu Meet Assoc Comput Linguist (ACL), Vol 2. 2022;1--9.

\bibitem{wang2022ln}
Qi W, Ruan YP, Zuo Y, Li T.
Parameter-efficient tuning on layer normalization for pre-trained language models.
arXiv preprint arXiv:2211.08682. 2022.

\bibitem{giannou2023layernorm}
Giannou A, Rajput S, Papailiopoulos D.
The expressive power of tuning only the normalization layers.
In: Proc Conf Learn Theory (COLT). 2023.

\bibitem{valizadeh2024layernorm}
ValizadehAslani T, Liang H.
LayerNorm: A key component in parameter-efficient fine-tuning.
arXiv preprint arXiv:2403.20284. 2024.

\bibitem{hu2022lora}
Hu EJ, Shen Y, Wallis P, Allen-Zhu Z, Li Y, Wang S, et al.
LoRA: Low-rank adaptation of large language models.
In: Int Conf Learn Represent (ICLR). 2022.

\bibitem{abdullah2026tap}
Abdullah A, Ebert N, Wasenm{\"u}ller O.
TAP into the patch tokens: Leveraging vision foundation model features for AI-generated image detection.
In: Proc IEEE/CVF Conf Comput Vis Pattern Recognit Workshops (CVPRW). 2026. arXiv:2604.26772.

\bibitem{xu2026bpl}
Xu Y, He R, Yan Z, Dong Y, Han Z, Wang X, Yin Y.
BPL: Generalizable deepfake detection via bias-only pair-aware learning.
In: Proc Int Conf Mach Learn (ICML). 2026.

\bibitem{zhang2026layermask}
Zhang X, Weng W, Fu D, Chen B, Li Z, He Z, Fu Z.
Generalizable deepfake detection based on forgery-aware layer masking and multi-artifact subspace decomposition.
arXiv preprint arXiv:2601.01041. 2026.

\bibitem{liu2024mole}
Liu Z, Wang H, Kang Y, Wang S.
Mixture of low-rank experts for transferable AI-generated image detection.
arXiv preprint arXiv:2404.04883. 2024.

\bibitem{chen2026envsubspace}
Chen S, Jia H, Li C, Ma C, Gao Z, Chen S.
Environment-invariant subspace learning for generalizable deepfake detection.
arXiv preprint arXiv:2608.17700. 2026.

\bibitem{liao2026dfa}
Liao J, Wei Y, Bon RCC, Wang S, Chow KP, Lam KY.
Deepfake Forensics Adapter: A dual-stream network for generalizable deepfake detection.
In: Proc Int Conf Digital Forensics and Cyber Crime (ICDF2C). 2025. arXiv:2603.01450.

\bibitem{han2024fcg}
Han YH, Huang TM, Hua KL, Chen JC.
Towards more general video-based deepfake detection through facial component guided adaptation for foundation model.
arXiv preprint arXiv:2404.05583. 2024.

\bibitem{lin2025repdfd}
Lin K, Lin Y, Li W, Yao T, Li B.
Standing on the shoulders of giants: Reprogramming visual-language model for general deepfake detection.
In: Proc AAAI Conf Artif Intell (AAAI). 2025. arXiv:2409.02664.

\bibitem{jiang2026adaptprompt}
Jiang Y, Alam MT, Khan SA, Dang-Nguyen DT, Karray F.
AdaptPrompt: Parameter-efficient adaptation of VLMs for generalizable deepfake detection.
In: Deepfake Forensics Workshop, ACM Multimedia. 2026. arXiv:2512.17730.

\bibitem{yang2026separable}
Yang E, Wu B, Li Y.
Generalizable face forgery detection via separable prompt learning.
arXiv preprint arXiv:2604.17307. 2026.

\bibitem{wang2020alignment}
Wang T, Isola P.
Understanding contrastive representation learning through alignment and uniformity on the hypersphere.
In: Proc Int Conf Mach Learn (ICML). 2020.

\bibitem{wang2026shortcut}
Wang Y, Yang Y, Liu J, Zhu F, Wu L, Chua TS.
Suppressing forgery-specific shortcuts for generalizable deepfake detection.
arXiv preprint arXiv:2606.01843. 2026.

\bibitem{wortsman2022soups}
Wortsman M, Ilharco G, Gadre SY, Roelofs R, Gontijo-Lopes R, Morcos AS, et al.
Model soups: Averaging weights of multiple fine-tuned models improves accuracy without increasing inference time.
In: Proc Int Conf Mach Learn (ICML). 2022.

\bibitem{wortsman2022wiseft}
Wortsman M, Ilharco G, Kim JW, Li M, Kornblith S, Roelofs R, et al.
Robust fine-tuning of zero-shot models.
In: Proc IEEE/CVF Conf Comput Vis Pattern Recognit (CVPR). 2022.

\bibitem{shiohara2022sbi}
Shiohara K, Yamasaki T.
Detecting deepfakes with self-blended images.
In: Proc IEEE/CVF Conf Comput Vis Pattern Recognit (CVPR). 2022.

\bibitem{chen2022sladd}
Chen L, Zhang Y, Song Y, Liu L, Wang J.
Self-supervised learning of adversarial example: Towards good generalizations for deepfake detection.
In: Proc IEEE/CVF Conf Comput Vis Pattern Recognit (CVPR). 2022.

\bibitem{lin2024cdfa}
Lin Y, Song W, Li B, Li Y, Ni J, Chen H, Li Q.
Fake it till you make it: Curricular dynamic forgery augmentations towards general deepfake detection.
In: Proc Eur Conf Comput Vis (ECCV). 2024.

\bibitem{pellegrini2025design}
Pellegrini L, Pandolfini S, Maltoni D, Ferrara M, Prati M, Ramilli M.
Generalized design choices for deepfake detectors.
arXiv preprint arXiv:2511.21507. 2025.

\bibitem{park2025mold}
Park N, Kim K, Choe J, Shim H.
Rethinking the use of vision transformers for AI-generated image detection.
arXiv preprint arXiv:2512.04969. 2025.

\bibitem{delibasoglu2026limits}
Delibasoglu I.
Cross-domain generalization limits of vision foundation models in facial deepfake detection.
arXiv preprint arXiv:2605.24965. 2026.

\bibitem{namazifar2023bias}
Namazifar M, Hazarika D, Hakkani-T{\"u}r D.
Role of bias terms in dot-product attention.
In: Proc IEEE Int Conf Acoust Speech Signal Process (ICASSP). 2023.

\bibitem{cao2022recce}
Cao J, Ma C, Yao T, Chen S, Ding S, Yang X.
End-to-end reconstruction-classification learning for face forgery detection.
In: Proc IEEE/CVF Conf Comput Vis Pattern Recognit (CVPR). 2022.

\bibitem{cheng2024prodet}
Cheng J, Yan Z, Zhang Y, Luo Y, Wang Z, Li C.
Can we leave deepfake data behind in training deepfake detector?
Adv Neural Inf Process Syst (NeurIPS). 2024.

\bibitem{huang2023iid}
Huang B, Wang Z, Yang J, Ai J, Zou Q, Wang Q, Ye D.
Implicit identity driven deepfake face swapping detection.
In: Proc IEEE/CVF Conf Comput Vis Pattern Recognit (CVPR). 2023.

\bibitem{liu2021spsl}
Liu H, Li X, Zhou W, Chen Y, He Y, Xue H, et al.
Spatial-phase shallow learning: Rethinking face forgery detection in frequency domain.
In: Proc IEEE/CVF Conf Comput Vis Pattern Recognit (CVPR). 2021.

\bibitem{ni2022core}
Ni Y, Meng D, Yu C, Quan C, Ren D, Zhao Y.
CORE: Consistent representation learning for face forgery detection.
In: Proc IEEE/CVF Conf Comput Vis Pattern Recognit Workshops (CVPRW). 2022.

\bibitem{yan2023ucf}
Yan Z, Zhang Y, Fan Y, Wu B.
UCF: Uncovering common features for generalizable deepfake detection.
In: Proc IEEE/CVF Int Conf Comput Vis (ICCV). 2023.

\bibitem{yan2024latent}
Yan Z, Luo Y, Lyu S, Liu Q, Wu B.
Transcending forgery specificity with latent space augmentation for generalizable deepfake detection.
In: Proc IEEE/CVF Conf Comput Vis Pattern Recognit (CVPR). 2024.

\end{thebibliography}

\end{document}
